\documentclass[12pt]{article}
\usepackage{amsmath, amssymb}
\usepackage[margin=1in]{geometry}
\setlength{\parskip}{6pt}
\title{QUBO Formulation: Remote-Sensing Image Segmentation}
\date{}
\begin{document}
\maketitle

\subsection*{Remote-Sensing Image Segmentation}

\paragraph{Problem Statement.}
Given a remote-sensing image represented as a grid graph with vertex set $V$ and edge set $E$, the task is to partition the pixels into two segments. Each pixel is assigned a label from a binary set, and the assignment must minimize a global segmentation energy. This energy comprises unary terms that encode the preference for assigning a specific label to a pixel, and pairwise terms that penalize boundary cuts between neighboring pixels with different labels. The goal is to determine the optimal binary labeling that minimizes this energy function.

\paragraph{Decision Variables.}
Let $N = |V|$ denote the total number of pixels in the image. For each pixel $i \in V$, we introduce a binary decision variable $x_i \in \{0,1\}$. The variable $x_i$ encodes the segment assignment for pixel $i$, where $x_i = 0$ corresponds to the first segment and $x_i = 1$ corresponds to the second segment. The complete solution is represented by the vector $x = (x_0, x_1, \dots, x_{N-1})^\top \in \{0,1\}^N$.

\paragraph{QUBO Derivation.}
The segmentation energy is directly expressed as a quadratic unconstrained binary optimization (QUBO) objective. We derive the corresponding Hamiltonian $H(x)$ and its coefficient matrix $Q$ through the following steps.

\textbf{Step 1.} The original objective function to be minimized is given by the sum of unary label costs and pairwise boundary penalties:
\begin{align}
  \min_{x \in \{0,1\}^N} \quad H(x) = \sum_{i \in V} a_i x_i + \sum_{(i,j) \in E} b_{ij} x_i x_j,
\end{align}
where $a_i$ denotes the unary cost coefficient for pixel $i$, and $b_{ij}$ denotes the pairwise boundary penalty coefficient for the edge connecting pixels $i$ and $j$.

\textbf{Step 2.} The problem formulation imposes no additional equality or inequality constraints beyond the binary domain restriction $x_i \in \{0,1\}$ for all $i \in V$. Therefore, no constraint equations need to be written.

\textbf{Step 3.} Since there are no constraints, no penalty terms are required. The formulation relies solely on the binary domain $\{0,1\}^N$, and the objective function is already in quadratic form. Consequently, the penalty contribution to the Hamiltonian is identically zero.

\textbf{Step 4.} Combining the objective and the absence of penalty terms, the single QUBO Hamiltonian $H(x)$ is:
\begin{align}
  H(x) = \sum_{i \in V} a_i x_i + \sum_{(i,j) \in E} b_{ij} x_i x_j.
\end{align}

\textbf{Step 5.} We read off the entries of the upper-triangular QUBO matrix $Q$ and the constant offset $c$ from the expanded Hamiltonian. For any $i, j \in V$, the matrix entries and offset are given in closed form as:
\begin{align}
  Q_{i,i} &= a_i, \\
  Q_{i,j} &= b_{ij} \quad \text{for } (i,j) \in E \text{ with } i < j, \\
  Q_{i,j} &= 0 \quad \text{for } (i,j) \notin E \text{ with } i < j, \\
  c &= 0.
\end{align}
The optimization problem is thus equivalent to $\min_{x \in \{0,1\}^N} x^\top Q x + c$.

\paragraph{Penalty Weight Justification.}
In constrained QUBO formulations, a penalty weight $\lambda$ is typically introduced to enforce hard constraints by adding squared violation terms to the objective. For this remote-sensing segmentation problem, the specification contains no constraints beyond the intrinsic binary domain of the variables. Consequently, no penalty weight $\lambda$ is defined or required. The formulation is inherently unconstrained, and the binary restriction $x_i \in \{0,1\}$ is naturally satisfied by the hardware or solver domain without auxiliary penalty terms.

\paragraph{Correctness Argument.}
Minimizing $H(x)$ over the domain $\{0,1\}^N$ recovers the optimal segmentation because the Hamiltonian is constructed to be exactly equivalent to the physical segmentation energy. Every binary assignment $x \in \{0,1\}^N$ corresponds to a valid segmentation of the image pixels, meaning all feasible solutions are evaluated directly without penalty. Since the objective function is a faithful representation of the segmentation cost, the global minimum of $H(x)$ over the binary hypercube must correspond to the labeling that achieves the minimum possible segmentation energy. Thus, any solver returning the ground state of $H(x)$ yields the optimal partition and its associated minimum energy value.

\paragraph{Illustrative Example.}
Consider a concrete instance with $N = 3$ pixels, where the pixel graph is fully connected such that $E = \{(0,1), (1,2), (0,2)\}$. The instance parameters are the unary costs $a_0, a_1, a_2$ and the pairwise weights $b_{01}, b_{12}, b_{02}$. Instantiating the general formulas from the derivation yields the following concrete objective:
\begin{align}
  H(x) &= a_0 x_0 + a_1 x_1 + a_2 x_2 + b_{01} x_0 x_1 + b_{12} x_1 x_2 + b_{02} x_0 x_2.
\end{align}
There are no penalty terms, as the problem remains unconstrained. The resulting QUBO matrix entries and offset are explicitly:
\begin{align}
  Q_{0,0} &= a_0, & Q_{1,1} &= a_1, & Q_{2,2} &= a_2, \\
  Q_{0,1} &= b_{01}, & Q_{1,2} &= b_{12}, & Q_{0,2} &= b_{02}, \\
  c &= 0.
\end{align}
For this $N=3$ instance, the binary search space contains $2^3 = 8$ possible assignments. The optimal solution $x^*$ is the binary vector in $\{0,1\}^3$ that minimizes the expression above, and the minimum energy is $H(x^*)$. Substituting numerical values for $a_i$ and $b_{ij}$ allows direct enumeration or quantum/classical solver execution to identify $x^*$ and $H(x^*)$.

\end{document}